% Two planar fields and a material element transported by their flow.
% Native size: 11.52 cm wide, with 9 pt labels in the 11 pt notes.
\begin{tikzpicture}[
  x=1cm,
  y=1cm,
  font=\footnotesize,
  field arrow/.style={
    draw=black,
    line width=0.5pt,
    line cap=round,
    -{Latex[length=1.4mm,width=1.05mm]}
  },
  initial element/.style={
    draw=black,
    line width=0.5pt,
    dash pattern=on 2pt off 1.8pt,
    line cap=butt,
    line join=miter
  },
  final element/.style={
    draw=black,
    line width=0.6pt,
    line join=miter
  },
  panel title/.style={font=\small,inner sep=0pt},
  field label/.style={inner sep=0pt},
  result label/.style={inner sep=0pt}
]
  \path[use as bounding box] (0,0) rectangle (11.52,6.0);

  % Preserve the examples and the flow interval of the original figure.
  % PGF uses degrees for angles; the rigid rotation is 0.32 radians.
  \def\dcFlowTime{0.32}
  \def\dcHalfSize{0.23}
  \pgfmathsetmacro{\dcExpandedHalfSize}{\dcHalfSize*exp(\dcFlowTime)}
  \pgfmathsetmacro{\dcRotationAngle}{\dcFlowTime*180/pi}

  % (a) F=(x,y,0): uniform local expansion, with no rotation.
  \begin{scope}[shift={(2.60,3.0)}]
    \node[panel title] at (0,2.70) {(a) Divergence};
    \node[field label] at (0,2.25) {$\mathbf F=(x,y,0)$};

    \begin{scope}[x=2cm,y=2cm]
      % Centred arrows, with length proportional to |F| (scale 1/2.8).
      \foreach \dcX in {-0.82,-0.41,0,0.41,0.82}{
        \foreach \dcY in {-0.82,-0.41,0,0.41,0.82}{
          \pgfmathtruncatemacro{\dcNonzero}{abs(\dcX)+abs(\dcY)>0}
          \ifnum\dcNonzero=1\relax
            \draw[field arrow]
              ({\dcX-\dcX/5.6},{\dcY-\dcY/5.6}) --
              ({\dcX+\dcX/5.6},{\dcY+\dcY/5.6});
          \fi
        }
      }
      \draw[initial element]
        (-\dcHalfSize,-\dcHalfSize) rectangle (\dcHalfSize,\dcHalfSize);
      \draw[final element]
        (-\dcExpandedHalfSize,-\dcExpandedHalfSize)
        rectangle (\dcExpandedHalfSize,\dcExpandedHalfSize);
      \fill[black] (0,0) circle[radius=0.8pt];
    \end{scope}

    \node[result label] at (0,-2.35) {$\nabla\!\cdot\!\mathbf F=2$};
    \node[result label] at (0,-2.75) {$\nabla\!\times\!\mathbf F=\mathbf 0$};
  \end{scope}

  % (b) F=(-y,x,0): anticlockwise rigid rotation, preserving area.
  \begin{scope}[shift={(8.92,3.0)}]
    \node[panel title] at (0,2.70) {(b) Curl};
    \node[field label] at (0,2.25) {$\mathbf F=(-y,x,0)$};

    \begin{scope}[x=2cm,y=2cm]
      \foreach \dcX in {-0.82,-0.41,0,0.41,0.82}{
        \foreach \dcY in {-0.82,-0.41,0,0.41,0.82}{
          \pgfmathtruncatemacro{\dcNonzero}{abs(\dcX)+abs(\dcY)>0}
          \ifnum\dcNonzero=1\relax
            \draw[field arrow]
              ({\dcX+\dcY/5.6},{\dcY-\dcX/5.6}) --
              ({\dcX-\dcY/5.6},{\dcY+\dcX/5.6});
          \fi
        }
      }
      \draw[initial element]
        (-\dcHalfSize,-\dcHalfSize) rectangle (\dcHalfSize,\dcHalfSize);
      \begin{scope}[rotate=\dcRotationAngle]
        \draw[final element]
          (-\dcHalfSize,-\dcHalfSize) rectangle (\dcHalfSize,\dcHalfSize);
      \end{scope}
      \fill[black] (0,0) circle[radius=0.8pt];
    \end{scope}

    \node[result label] at (0,-2.35) {$\nabla\!\cdot\!\mathbf F=0$};
    \node[result label] at (0,-2.75) {$\nabla\!\times\!\mathbf F=2\hat{\mathbf z}$};
  \end{scope}
\end{tikzpicture}
